\chapter{Electroconvulsive Therapy}

\section*{What Is This Test or Procedure?}
Electroconvulsive therapy (ECT) is a treatment that delivers a controlled electrical stimulus to the brain to induce a brief, generalized seizure, used for severe psychiatric disorders.

\section*{Alternative Names}
ECT; shock therapy; electroshock treatment

\section*{Principle}
Under anesthesia, a small electrical current is applied through scalp electrodes to provoke a therapeutic seizure. Repeated treatments produce changes in brain circuits and neurochemistry that relieve severe depression and other conditions.

\section*{History}
ECT was introduced by Cerletti and Bini in 1938 and, despite stigma, remains among the most effective treatments for severe depression.

\section*{Why This Test or Procedure Is Ordered}
Performed for severe, psychotic, or catatonic depression, acute mania, and suicidality, especially when medications fail or rapid response is needed.

\section*{Is This a Primary or Secondary Test or Procedure}
Primary treatment for severe, treatment-resistant psychiatric conditions.

\section*{How This Test or Procedure Is Performed}
Performed under general anesthesia with muscle relaxants and EEG and ECG monitoring. A brief electrical pulse induces a seizure lasting under a minute. A course typically includes 6-12 sessions.

\section*{What Preparation Is Needed}
Pre-anesthesia evaluation, discontinuation of interacting medications where appropriate, and fasting before each session.

\section*{Contraindications and Precautions}
Recent myocardial infarction, stroke, or intracranial mass with increased pressure are relative contraindications. Precaution: anesthesia and monitoring are mandatory.

\section*{Risks}
Transient memory loss and confusion, headache, nausea, and cardiovascular changes.

\section*{What Happens After the Test or Procedure}
The patient recovers from anesthesia and is monitored; memory side effects usually resolve over days to weeks.

\section*{Understanding Your Results}
Response is judged by improvement in mood, psychomotor status, and safety over the treatment course.

\section*{Normal Results}
Substantial improvement in depressive or psychotic symptoms during the treatment course.

\section*{Follow-up Tests and Procedures}
Maintenance ECT or medication adjustments prevent relapse.

\section*{References}
\begin{enumerate}
  \item MedlinePlus. Electroconvulsive therapy. U.S. National Library of Medicine; 2025.
  \item Current Medical Diagnosis and Treatment. McGraw-Hill; 2025.
\end{enumerate}

\clearpage
